% 源分章 PDF 第 1—6 页；章首导言与第 12.1 节。
\chapter{样本平均处理效应}
\begingroup
\raggedbottom
\chapterauthors{Laura B. Balzer、Maya L. Petersen 与 Mark J. van der Laan}
\chapteraffiliations{%
L. B. Balzer（通讯作者）：美国马萨诸塞大学阿默斯特分校公共卫生与健康科学学院生物统计学与流行病学系，416 Arnold House, 715 North Pleasant St, Amherst, MA 01003, USA；电子邮箱：\href{mailto:lbalzer@umass.edu}{lbalzer@umass.edu}。\par
M. L. Petersen：美国加利福尼亚大学伯克利分校流行病学部及生物统计学部，101 Haviland Hall, \#7358, Berkeley, CA 94720, USA；电子邮箱：\href{mailto:mayaliv@berkeley.edu}{mayaliv@berkeley.edu}。\par
M. J. van der Laan：美国加利福尼亚大学伯克利分校生物统计学部及统计学系，101 Haviland Hall, \#7358, Berkeley, CA 94720, USA；电子邮箱：\href{mailto:laan@berkeley.edu}{laan@berkeley.edu}。\par
本章 DOI：\url{https://doi.org/10.1007/978-3-319-65304-4_12}。}

在整群随机试验（cluster randomized trial，CRT）中，研究单位通常不是从某个界定清晰的目标总体中抽取的简单随机样本。相反，目标总体往往是假想的，或界定不清；研究单位的选择则往往具有系统性，受后勤和实际条件的考虑所驱动。因此，总体平均处理效应（population average treatment effect，PATE）可能既没有明确的定义，也不容易解释。相比之下，样本平均处理效应（sample average treatment effect，SATE）是研究单位反事实结局之差的均值。当研究单位并非抽取自某个特定的、受关注的超总体（super-population）时，样本参数容易解释，而且可以说是最相关的参数。此外，在多数情形下，对样本参数的估计比对总体参数的估计更有效。

\begin{sourcebox}
本章展示如何在采用和未采用配对匹配（pair-matching）的试验中，利用靶向最小损失估计（targeted minimum loss-based estimation，TMLE）对样本效应进行估计和推断。我们研究样本效应 TMLE 的渐近性质和有限样本性质，并给出一个保守方差估计量（conservative variance estimator）。有限样本模拟说明，将样本效应选为推断目标，可能提高精度和统计功效（statistical power）。本章改编自 Balzer 等（2016c）。
\end{sourcebox}

许多研究旨在估计暴露对关注结局的影响。目标因果参数往往是 PATE：某个总体全部成员均暴露时的反事实结局，与该总体全部成员均未暴露时的反事实结局之差的期望。如果没有未测量混杂因素，且暴露分配有充分的变异，即随机化假设和正值性假设成立，那么就能将 PATE 识别为观测数据分布的函数（Rosenbaum 和 Rubin，1983b；Robins，1986）。所得统计参数可以用多种算法估计，包括匹配和逆加权估计量（例如 Horvitz 和 Thompson，1952；Rosenbaum 和 Rubin，1983b；Shen 等，2014）、简单代入估计量（例如 Robins，1986；Snowden 等，2011），以及双重稳健算法（例如 Robins 等，1994；van der Laan 和 Robins，2003；van der Laan 和 Rubin，2006；van der Laan 和 Rose，2011）。

另一个因果参数是 SATE（Neyman，1923；Rubin，1990；Imbens，2004；Imai，2008；Schochet，2013；Imbens 和 Rubin，2015）。样本效应是实际研究单位反事实结局之差的平均。将 SATE 选为关注参数具有若干潜在优势。首先，SATE 很容易解释为手头样本的干预效应。其次，SATE 避免了关于从某个“界定模糊的研究单位超总体”中随机抽样、并将结果推广至该总体的假设（Schochet，2013）。换言之，即使这些单位是以系统性方式选入研究的，样本参数仍具有相关性和可解释性；这种选样方式在 CRT 中可能很常见。将研究结果扩展到更广泛或不同的总体，可以作为一个独立的研究问题，借助正式工具处理（例如 Cole 和 Stuart，2010；Stuart 等，2011；Bareinboim 和 Pearl，2013；Hartman 等，2015），而不必在指定参数时就假定这种推广成立。最后，同一个估计量用于估计样本效应时，通常比用于估计总体效应时具有更高的精度（Neyman，1923；Rubin，1990；Imbens，2004；Imai，2008）。

对于随机试验，Neyman（1923）最早提出用未调整估计量估计 SATE，即处理单位的结局均值与对照单位的结局均值之差。在这种情形下，给定研究单位的全部反事实结局，均值差估计量对 SATE 无偏。但其方差仍不可识别，因为它依赖于反事实结局之间的相关性（Neyman，1923；Rubin，1990；Imbens，2004；Imai，2008）。后来，Imbens（2004）将这一研究推广到观察性研究中的有效估计量，即影响曲线等于有效影响曲线的正则渐近线性估计量。具体而言，他证明，给定研究单位的基线协变量及反事实结局，总体效应的有效估计量对样本效应无偏。他还将 SATE 有效估计量的方差表示为：同一个估计量用于 PATE 时的方差，减去总体中单位特定处理效应的方差。这表明，除非处理效应没有变异，否则标准方差估计量会存在向上偏差。

\begin{sourcebox}
本章介绍一种 TMLE，用于在采用和未采用配对匹配的试验中估计样本效应并进行推断。我们的结果推广了 Imbens（2004）的方差推导，允许结局回归（即给定暴露和协变量的条件结局均值）误设、对倾向评分（即给定协变量后接受暴露的条件概率）进行估计，以及自适应配对匹配（Balzer 等，2015）。配对匹配是 CRT 中常用的一种设计策略，可以保障研究的可信度并提高功效（Klar 和 Donner，1997；Greevy 等，2004；Imai 等，2009；Hayes 和 Moulton，2009；van der Laan 等，2013a；Balzer 等，2015）。据我们所知，TMLE 是首个被提出用于配对匹配试验中样本效应的有效估计量。
\end{sourcebox}

作为引导性例子，考虑一项假想的 HIV 预防与治疗 CRT。假设试验包含 $n=30$ 个社区。干预社区定期提供 HIV 检测，所有检测为 HIV 阳性者均立即具有接受抗逆转录病毒治疗（antiretroviral therapy，ART）的资格，同时获得加强的治疗启动、依从性和持续随诊服务。对照社区则按照现行标准诊疗，为所有 HIV 检测阳性者提供 ART。主要假设是：普遍实施“检测并治疗”（test-and-treat）策略，将降低试验期间 HIV 的累积发生率。为便于讨论，我们着重考虑社区层面的数据。因此，这些结果同样适用于有聚集性和无聚集性的数据结构。

\section{因果模型与因果参数}

考虑一项两组随机试验的如下数据生成过程。首先选择研究单位。有些试验从一个界定清晰的目标总体中获取简单随机样本；另一些研究则可能没有明确的目标总体，既不清楚单位抽样自何处，也不清楚希望对哪个总体进行推断。例如，在 SEARCH 试验中，从乌干达西部的姆巴拉拉（Mbarara）地区、东部的托罗罗（Tororo）地区以及肯尼亚南尼安扎（Southern Nyanza）省选出 32 个社区。具体做法是：先对满足纳入标准的 54 个候选社区进行民族志调查与制图；纳入标准包括社区规模、医疗保健基础设施，以及是否有维护良好的交通路线可达等。随后，根据地区、职业构成、迁移指数等多项特征，选出匹配最好的 16 对社区（Balzer 等，2015）。选定研究单位后，通常还会测量其他协变量。在贯穿本章的例子中，另外收集的协变量可能包括男性包皮环切覆盖率、HIV 患病率指标，以及社区层面的 HIV RNA 病毒载量指标。下文始终以 $W$ 表示基线协变量。

接下来，将干预随机分配给研究单位。可以将 $n/2$ 个单位随机分配至干预组，其余单位分配至对照组，也可以在匹配对内随机化，从而保证干预的等额分配。为便于表述，我们针对干预完全随机分配这一简单情形介绍因果模型，但所得结果具有一般性；向配对匹配试验的扩展见第 12.4 节。令 $A$ 为表示所分配干预水平的二元变量。在假想的 CRT 中，社区被分配至处理组（普遍检测并治疗）时，$A$ 等于 1；被分配至对照组（标准诊疗）时，$A$ 等于 0。在随访结束时测量结局 $Y$。在我们的试验中，$Y$ 为相关时间段内 HIV 的累积发生率。因此，某个研究单位的观测数据为
\[
O=(W,A,Y).
\]
假设观测到 $O$ 的 $n$ 个独立同分布副本，其分布为 $P_0$。回顾一下，下标 0 表示观测数据的真实分布。对于样本平均处理效应和条件平均处理效应的估计与推断，可以通过给定基线协变量向量 $(W_1,W_2,\ldots,W_n)$ 来弱化独立同分布假设（Balzer 等，2015）。

以下结构因果模型（structural causal model）描述了这一数据生成过程（Pearl，1995，2009a）。假设观测数据的每个分量，都是其父变量（可能影响其取值的变量）和不可观测背景因素的确定性函数：
\begin{equation}
\begin{aligned}
W&=f_W(U_W),\\
A&=\mathbb I(U_A<0.5),\\
Y&=f_Y(W,A,U_Y),
\end{aligned}
\tag{12.1}
\end{equation}
其中，背景因素集合 $U=(U_W,U_A,U_Y)$ 具有某个联合分布 $P_U$。按照试验设计，决定干预分配的随机误差 $U_A$，独立于影响基线协变量的未测量因素 $U_W$ 和影响结局的未测量因素 $U_Y$：
\[
U_A\indep(U_W,U_Y).
\]
具体而言，$U_A$ 独立地抽取自 $\operatorname{Uniform}(0,1)$。这一因果模型 $\mathcal M^{\mathcal F}$ 蕴含了观测数据 $O$ 的所有可能分布所构成的统计模型。在随机试验中，统计模型 $\mathcal M$ 是半参数模型。

通过对结构因果模型施加干预，可以生成反事实结局 $Y_a$，即将该单位分配至 $A=a$ 时的结局，这种分配可能与事实相反：
\[
\begin{aligned}
W&=f_W(U_W),\\
A&=a,\\
Y_a&=f_Y(W,a,U_Y).
\end{aligned}
\]
在这一框架中，反事实结局 $Y_a$ 是随机变量。在贯穿本章的例子中，$Y_a$ 是将社区分配至处理水平 $A=a$ 时的反事实 HIV 累积发生率，这种分配可能与事实相反。

随后，可以利用反事实的分布定义所关注的因果参数。推断目标往往是总体平均处理效应：
\[
\mathrm{PATE}=\EE[Y_1-Y_0].
\]
它是这些单位所抽样自的潜在目标总体中，反事实结局之差的期望。从结构因果模型可以看出，期望是对测量因素 $W$ 和未测量因素 $U_Y$ 取的，二者共同决定总体的反事实结局。换言之，PATE 的真值不依赖于抽样得到的 $W$ 或 $U_Y$ 的取值。在假想试验中，PATE 是以下两个反事实 HIV 累积发生率期望之差：某个假想目标总体的全部社区均实施检测并治疗策略时的期望，与该假想目标总体的全部社区均继续标准诊疗时的期望；这两种情形都可能与事实相反。

另一种因果参数是样本平均处理效应，最早由 Neyman（1923）提出：
\[
\mathrm{SATE}=\frac1n\sum_{i=1}^n[Y_{1,i}-Y_{0,i}].
\]
它就是这 $n$ 个研究单位的干预效应。SATE 是一个数据自适应参数，其取值依赖于纳入研究的单位。有关其他数据自适应参数估计与推断的近期研究，参见第 9 章。即使不存在研究单位所选自的明确超总体，SATE 仍具有可解释性。在我们的例子中，SATE 是这 $n$ 个研究社区在检测并治疗策略与标准诊疗下，反事实 HIV 累积发生率之差的平均。

在 CRT 中，靶向样本效应相对于靶向总体效应可能具有若干优势。首先，研究单位可能并非抽样自某个唯一的真实世界目标总体，我们也可能并不希望对这样一个总体进行推断；这里所说的真实世界总体区别于假想总体。虽然恰当的分析方法能够减轻对系统性选样的担忧，但所得 PATE 估计的解释及政策相关性仍可能不清楚。相比之下，靶向 SATE 能够严格估计一个界定清晰、由研究社区全体居民构成的真实世界总体中的干预效应。所得 SATE 估计不依赖于有关抽样机制的任何假设，解释明确，而且通常比 PATE 估计具有更高精度。如下文所述，样本效应估计量的功效至少与总体效应估计量一样高；存在效应修饰时，预计前者的功效更高（Rubin，1990；Imbens，2004；Imai，2008）。

当然，将 CRT 中发现的效应迁移至新的总体和情境，仍然具有重要政策意义。但是，其他真实世界目标总体很可能在多个方面不同于当前情境，而且这些差别很可能影响效应大小。因此，SATE 和 PATE 都不能直接适用于这些新情境。希望结果具有可推广性，并不能构成优先选择总体参数而不是样本参数的理由。较好的做法是，将干预效应向当前样本以外情境的推广（或迁移）作为一个独立研究问题，充分利用现有的现代工具（例如 Cole 和 Stuart，2010；Stuart 等，2011；Bareinboim 和 Pearl，2013；Hartman 等，2015）。
